\documentclass[../DoAn.tex]{subfiles}
\begin{document}
Chương~\ref{chapter:Introduction} đã trình bày bài toán chuyển hóa video bài giảng thụ động thành bài học chủ động cùng định hướng giải pháp tổng quát. Để có cơ sở thiết kế hệ thống, chương này tiến hành khảo sát hiện trạng và phân tích yêu cầu một cách chi tiết. Phần~\ref{section:2.1} khảo sát hiện trạng từ ba nguồn gồm nhu cầu người dùng, các hệ thống đã có và các ứng dụng tương tự, kèm phân tích so sánh ưu nhược điểm. Phần~\ref{section:2.2} trình bày tổng quan chức năng thông qua biểu đồ use case và quy trình nghiệp vụ. Phần~\ref{section:2.3} đặc tả chi tiết các chức năng quan trọng. Cuối cùng, Phần~\ref{section:2.4} nêu các yêu cầu phi chức năng mà hệ thống cần đáp ứng.

\section{Khảo sát hiện trạng}
\label{section:2.1}
Việc khảo sát hiện trạng và yêu cầu của phần mềm được thực hiện trên ba nguồn chính: nhu cầu của người dùng, các hệ thống đã có liên quan đến bài toán, và các ứng dụng tương tự đang được sử dụng phổ biến. Mục tiêu của khảo sát là làm rõ khoảng trống mà các sản phẩm hiện có chưa giải quyết, đồng thời xác định tập tính năng quan trọng cần phát triển.

\subsection{Nhu cầu người dùng}
\label{subsection:2.1.1}
Hệ thống hướng tới hai nhóm người dùng có nhu cầu khác biệt rõ rệt. Nhóm thứ nhất là người học, gồm sinh viên và những người tự học qua video. Như đã phân tích ở Chương~\ref{chapter:Introduction}, hình thức học một chiều khiến người học dễ thụ động, xem lướt và ngộ nhận đã hiểu bài trong khi chưa thể tái hiện hay vận dụng kiến thức~\cite{prince2004active,karpicke2008critical}. Từ những phân tích trên, hệ thống cần đáp ứng nhu cầu của người học là được kiểm tra hiểu bài thường xuyên ngay trong lúc xem, nhận phản hồi tức thời thay vì chờ đến bài kiểm tra cuối khóa, và được củng cố kiến thức ở những thời điểm hợp lý. Những nhu cầu này phù hợp với các nguyên lý đã được kiểm chứng trong khoa học học tập như học tập chủ động, hiệu ứng kiểm tra và việc thiết kế tương tác hiệu quả quanh video~\cite{brame2016effective}.

Nhóm thứ hai là giảng viên và người phụ trách tổ chức đào tạo. Nhóm này thường đã có sẵn một kho video bài giảng nhưng thiếu thời gian và công cụ để biến chúng thành bài học tương tác. Việc soạn câu hỏi xen kẽ video một cách thủ công đòi hỏi xem lại toàn bộ video, xác định các mốc thời gian thích hợp và nhập nội dung cho từng câu hỏi, một công việc tốn nhiều công sức và khó duy trì khi số lượng bài giảng tăng lên. Vì vậy, hệ thống cần cung cấp một công cụ tự động phân tích video và đề xuất sẵn các hoạt động tương tác, cho phép xem trước và chỉnh sửa kết quả trước khi đưa tới người học, đồng thời hỗ trợ tổ chức nội dung theo cấu trúc khóa học và quản lý người học tham gia.

\subsection{Các hệ thống đã có}
\label{subsection:2.1.2}
Trong môi trường giáo dục, các hệ thống quản lý học tập như Moodle\footnote{\url{https://moodle.org}} hay Google Classroom\footnote{\url{https://classroom.google.com}} đã được sử dụng rộng rãi để tổ chức khóa học, phân phối tài liệu và quản lý bài nộp. Tuy nhiên, các hệ thống này tập trung vào khâu quản lý và phân phối nội dung hơn là vào trải nghiệm tương tác bên trong từng video. Video thường được nhúng như một tài nguyên tĩnh, còn việc đánh giá hiểu bài được tách rời thành các bài trắc nghiệm hoặc bài tập riêng đặt sau video. Hệ quả là khâu kiểm tra hiểu bài bị ngắt khỏi dòng chảy của bài giảng, làm giảm tính tức thời của phản hồi và không tận dụng được thời điểm phù hợp để củng cố kiến thức ngay khi nó vừa được trình bày. Ngoài ra, toàn bộ nội dung tương tác trong các hệ thống này vẫn phải do giảng viên tạo thủ công.

\subsection{Các ứng dụng tương tự}
\label{subsection:2.1.3}
Đồ án khảo sát ba nhóm ứng dụng tiêu biểu có liên quan trực tiếp đến bài toán: các nền tảng khóa học trực tuyến mở đại trà, các công cụ tạo video tương tác chuyên dụng, và các kho video bài giảng công cộng.

Nhóm thứ nhất là các nền tảng khóa học trực tuyến mở đại trà (MOOC) như Coursera\footnote{\url{https://www.coursera.org}} và edX\footnote{\url{https://www.edx.org}}. Các nền tảng này có quy mô lớn, chất lượng khóa học cao cùng lộ trình học bài bản và chứng chỉ (Hình~\ref{fig:khaosat_coursera}). Tuy nhiên, hoạt động đánh giá chủ yếu là các bài trắc nghiệm hoặc bài tập đặt ở cuối bài học hoặc cuối tuần và đều do giảng viên biên soạn thủ công; mức độ tương tác ngay trong lúc xem video còn hạn chế. Đặc điểm một chiều này là một trong những nguyên nhân dẫn đến tỷ lệ hoàn thành thấp đã đề cập ở Chương~\ref{chapter:Introduction}.

\begin{figure}[H]
\centering
\includegraphics[width=0.8\textwidth]{Hinhve/khaosat_coursera.png}
\caption{Giao diện học một khóa học trên nền tảng Coursera (nguồn: Coursera)}
\label{fig:khaosat_coursera}
\end{figure}

Nhóm thứ hai là các công cụ tạo video tương tác chuyên dụng như H5P\footnote{\url{https://h5p.org}} và Edpuzzle\footnote{\url{https://edpuzzle.com}}. Các công cụ này cho phép giảng viên chèn nhiều dạng câu hỏi (trắc nghiệm, đúng/sai, điền khuyết, tự luận) vào các mốc thời gian bất kỳ trong video; người học phải trả lời để tiếp tục xem và nhận phản hồi ngay lập tức, nhờ đó tính tương tác trong video được cải thiện rõ rệt so với nhóm MOOC. Gần đây, một số công cụ như Edpuzzle đã bổ sung tính năng gợi ý câu hỏi bằng trí tuệ nhân tạo, song mới dừng ở mức hỗ trợ gợi ý và người dạy vẫn phải tự kiểm duyệt, chỉnh sửa. Nhìn chung, hạn chế cốt lõi vẫn còn: giảng viên phải tự xác định vị trí chèn và biên tập nội dung cho từng câu hỏi (Hình~\ref{fig:khaosat_edpuzzle}), nên gánh nặng soạn thảo gần như không được giảm bớt.

\begin{figure}[H]
\centering
\includegraphics[width=0.8\textwidth]{Hinhve/khaosat_edpuzzle.png}
\caption{Chèn câu hỏi vào các mốc thời gian của video bằng Edpuzzle (nguồn: Edpuzzle)}
\label{fig:khaosat_edpuzzle}
\end{figure}

Nhóm thứ ba là các kho video bài giảng công cộng như YouTube\footnote{\url{https://www.youtube.com}}, nổi bật về khả năng tiếp cận và khối lượng nội dung khổng lồ. Tuy nhiên, các nền tảng này gần như không tích hợp cơ chế tương tác hay đánh giá hiểu bài; người học chỉ xem một chiều và tự đánh giá mức độ hiểu của mình (Hình~\ref{fig:khaosat_youtube}).

\begin{figure}[H]
\centering
\includegraphics[width=0.8\textwidth]{Hinhve/khaosat_youtube.png}
\caption{Giao diện xem video bài giảng trên YouTube (nguồn: YouTube)}
\label{fig:khaosat_youtube}
\end{figure}

Bảng~\ref{table:sosanh} so sánh ba nhóm sản phẩm trên cùng với hệ thống Dyadia đề xuất theo sáu tiêu chí phản ánh trực tiếp các nhu cầu đã phân tích: khả năng tương tác ngay trong video, mức độ tự động hóa khâu sinh nội dung, năng lực chấm điểm và phản hồi tự động, sự đa dạng của loại hình tương tác, công sức soạn thảo mà giảng viên phải bỏ ra, và khả năng phản hồi tức thời cho người học. Riêng hệ thống Dyadia hỗ trợ nhiều loại hoạt động tương tác, gồm câu hỏi trắc nghiệm, điền khuyết, bài nghe, bài đọc và bài viết tự luận; chi tiết được trình bày trong Chương~\ref{chapter:Experiment}.

\begin{table}[H]
\centering
\begin{tabular}{|p{3.6cm}|c|c|c|c|}
\hline
\textbf{Tiêu chí} & \textbf{MOOC} & \textbf{H5P/Edpuzzle} & \textbf{YouTube} & \textbf{Dyadia} \\ \hline
Tương tác trong video & Không & Có & Không & Có \\ \hline
Tự động sinh nội dung & Không & Một phần & Không & Có \\ \hline
Tự động chấm và phản hồi & Một phần & Một phần & Không & Có \\ \hline
Đa dạng loại tương tác & Thấp & Cao & Không & Cao \\ \hline
Công sức soạn thảo & Cao & Cao & Không áp dụng & Thấp \\ \hline
Phản hồi tức thời & Một phần & Có & Không & Có \\ \hline
\end{tabular}
\caption{So sánh các nhóm sản phẩm hiện có với hệ thống đề xuất}
\label{table:sosanh}
\end{table}

Phân tích Bảng~\ref{table:sosanh} cho thấy mỗi nhóm sản phẩm đều có điểm mạnh riêng nhưng vẫn tồn tại hạn chế cốt lõi. Các nền tảng MOOC mạnh về tổ chức khóa học nhưng tương tác trong video yếu và đòi hỏi nhiều công soạn thảo. Các công cụ như H5P và Edpuzzle giải quyết tốt việc tương tác trong video và phản hồi tức thời, nhưng khả năng tự động sinh nội dung mới dừng ở mức gợi ý từng câu hỏi rời rạc chứ chưa tạo trực tiếp một bài học hoàn chỉnh từ video, nên người dạy vẫn phải tự đặt và biên tập phần lớn nội dung. YouTube vượt trội về khả năng tiếp cận nhưng gần như không hỗ trợ tương tác hay đánh giá. Giá trị ``Một phần'' ở cột MOOC phản ánh việc một số nền tảng có bài kiểm tra tích hợp nhưng đặt ngoài luồng phát video, không phải trong khi xem. Điểm chung quan trọng nhất là trong phạm vi khảo sát, chưa nhóm sản phẩm nào kết hợp được đồng thời ba yếu tố: tự động sinh nội dung tương tác trực tiếp từ nội dung video, chấm điểm và phản hồi tự động, và đa dạng loại hình tương tác. Đây chính là khoảng trống mà hệ thống Dyadia hướng tới giải quyết.

Từ kết quả khảo sát, đồ án xác định các tính năng quan trọng cần phát triển gồm: (i) quản lý tổ chức, khóa học, chương, bài giảng và thành viên khóa học, bao gồm duyệt yêu cầu tham gia; (ii) tải lên, lưu trữ và phát video bài giảng; (iii) tự động phân tích video và sinh các hoạt động tương tác; (iv) thực hiện bài học có tương tác kèm chấm điểm, phản hồi tự động và theo dõi kết quả học tập; (v) theo dõi trạng thái và xử lý lỗi cho các tác vụ phân tích video và sinh nội dung tương tác; và (vi) quản lý người dùng cùng phân quyền theo vai trò. Các tính năng này được trình bày tổng quan trong Phần~\ref{section:2.2}.

\section{Tổng quan chức năng}
\label{section:2.2}

\subsection{Biểu đồ use case tổng quát}
\label{subsection:2.2.1}
Trên cơ sở tập tính năng đã xác định ở Phần~\ref{section:2.1}, phần này trình bày tổng quan chức năng của hệ thống Dyadia bằng biểu đồ use case. Hệ thống có bốn tác nhân cụ thể là Khách, Người học, Giảng viên và Quản trị viên, cùng một tác nhân tổng quát Người dùng đóng vai trò lớp cha cho các vai trò đã đăng nhập; biểu đồ use case tổng quát ở Hình~\ref{fig:usecase_tongquat} thể hiện các tác nhân này cùng những nhóm use case chính mà mỗi tác nhân có thể thực hiện.

\begin{figure}[H]
\centering
\includegraphics[width=\textwidth]{Hinhve/usecase_tongquat.png}
\caption{Biểu đồ use case tổng quát của hệ thống Dyadia}
\label{fig:usecase_tongquat}
\end{figure}

\textbf{a) Các tác nhân tham gia}

\textbf{Khách}: Người dùng tương tác với hệ thống khi chưa đăng nhập; tác nhân này chỉ có thể đăng ký tài khoản và đăng nhập.

\textbf{Người dùng}: Tác nhân tổng quát đại diện cho người dùng đã đăng nhập vào hệ thống bằng tên đăng nhập và mật khẩu. Đây là lớp cha của ba tác nhân Người học, Giảng viên và Quản trị viên; cả ba kế thừa từ tác nhân này các chức năng tài khoản dùng chung gồm cập nhật hồ sơ cá nhân, đổi mật khẩu và đăng xuất.

\textbf{Người học}: Người dùng đã đăng nhập với vai trò học viên, là sinh viên hoặc người tự học; tham gia khóa học, thực hiện các bài giảng có tương tác và theo dõi kết quả của bản thân.

\textbf{Giảng viên}: Người dùng đã đăng nhập với vai trò phụ trách nội dung trong một tổ chức đào tạo; tổ chức khóa học, đưa video bài giảng vào hệ thống, biên tập nội dung tương tác và theo dõi kết quả học tập của học viên.

\textbf{Quản trị viên}: Người dùng đã đăng nhập với vai trò quản trị toàn hệ thống; quản lý người dùng, tổ chức, phân quyền và giám sát các tác vụ phân tích video.

\textbf{b) Mô tả các use case chính}

Use case \textbf{Tham gia khóa học}: Người học gửi yêu cầu tham gia một khóa học và chờ giảng viên duyệt; sau khi được chấp nhận, người học truy cập được các bài giảng của khóa học đó.

Use case \textbf{Học bài giảng tương tác}: Người học vừa xem video bài giảng vừa trả lời các hoạt động tương tác (trắc nghiệm, điền khuyết, bài nghe, bài đọc, bài viết tự luận) được chèn xen kẽ tại các mốc thời gian của video. Use case này bao hàm (\(\ll\)include\(\gg\)) việc chấm điểm và phản hồi tự động ngay sau mỗi câu trả lời; kết quả từng bài được lưu lại.

Use case \textbf{Xem kết quả và phản hồi}: Người học xem lại điểm số và đáp án kèm nhận xét của các bài đã làm, đồng thời theo dõi tiến độ học tập của bản thân trên toàn khóa học.

Use case \textbf{Quản lý khóa học, chương và bài giảng}: Giảng viên tạo và tổ chức nội dung giảng dạy theo cấu trúc ba cấp khóa học -- chương -- bài giảng.

Use case \textbf{Tải lên video bài giảng}: Giảng viên tải tệp video lên hệ thống và gắn vào một bài giảng; video được lưu trữ và sẵn sàng cho bước phân tích sinh nội dung tương tác.

Use case \textbf{Phân tích video sinh nội dung tương tác}: Giảng viên kích hoạt quá trình phân tích tự động đối với video đã tải lên: hệ thống trích xuất lời thoại bằng công nghệ nhận dạng tiếng nói, phân đoạn nội dung, rồi dùng mô hình ngôn ngữ lớn để sinh các hoạt động tương tác cho từng đoạn. Do quá trình xử lý kéo dài, use case này bao hàm (\(\ll\)include\(\gg\)) việc theo dõi trạng thái tác vụ, cho phép giảng viên xem tiến độ và hủy tác vụ khi cần.

Use case \textbf{Biên tập bản chép lời}: Giảng viên chỉnh sửa bản chép lời mà hệ thống trích xuất tự động, gồm sửa nội dung lời thoại, gộp hoặc tách phân đoạn và điều chỉnh ranh giới các phân đoạn nội dung.

Use case \textbf{Biên tập hoạt động tương tác}: Giảng viên xem trước, chỉnh sửa, thêm mới hoặc xóa từng hoạt động tương tác mà hệ thống đã sinh trước khi đưa bài giảng tới người học. Khi chưa hài lòng với kết quả tự động, giảng viên có thể dùng phần mở rộng (\(\ll\)extend\(\gg\)) sinh lại hoạt động bằng trí tuệ nhân tạo.

Use case \textbf{Quản lý thành viên và duyệt yêu cầu tham gia}: Giảng viên xem danh sách thành viên khóa học, trực tiếp thêm hoặc loại bỏ học viên, và duyệt (chấp nhận hoặc từ chối) các yêu cầu tham gia đang chờ.

Use case \textbf{Theo dõi kết quả học tập}: Giảng viên xem các thống kê tổng hợp về mức độ tiếp thu của học viên, gồm tỷ lệ trả lời đúng theo từng câu hỏi, biểu đồ kết quả làm bài theo từng mốc thời gian của video và danh sách người học cần hỗ trợ.

Use case \textbf{Quản lý người dùng và phân quyền}: Quản trị viên tạo, khóa hoặc cập nhật tài khoản người dùng, gán vai trò và quyền truy cập trên toàn hệ thống.

Use case \textbf{Quản lý tổ chức và thành viên}: Quản trị viên tạo và cập nhật các tổ chức đào tạo, quản lý thành phần thành viên của từng tổ chức và gán vai trò cho các thành viên đó.

Use case \textbf{Giám sát tác vụ phân tích}: Quản trị viên theo dõi toàn bộ tác vụ phân tích video trên hệ thống, xem chi tiết trạng thái và thông tin lỗi của từng tác vụ, và hủy tác vụ khi cần.

Các use case có tính tổng hợp nêu trên được phân rã chi tiết trong các mục tiếp theo.

\subsection{Biểu đồ use case phân rã Tham gia khóa học}
\label{subsection:uc_thamgia}
Use case \emph{Tham gia khóa học} của người học được phân rã như Hình~\ref{fig:uc_thamgia}.

\begin{figure}[H]
\centering
\includegraphics[width=0.55\textwidth]{Hinhve/usecasephanra/usecasephanra_thamgiakhoahoc.png}
\caption{Biểu đồ use case phân rã Tham gia khóa học}
\label{fig:uc_thamgia}
\end{figure}

\emph{Tìm kiếm và xem khóa học} cho phép người học duyệt danh sách khóa học và xem thông tin chi tiết của một khóa học. \emph{Gửi yêu cầu tham gia} cho phép người học đăng ký vào khóa học; yêu cầu được chuyển tới giảng viên chờ duyệt. \emph{Xem trạng thái yêu cầu} cho phép người học theo dõi yêu cầu của mình đang chờ duyệt, đã được chấp nhận hay bị từ chối.

\subsection{Biểu đồ use case phân rã Học bài giảng tương tác}
\label{subsection:uc_hocbai}
Use case \emph{Học bài giảng tương tác} của người học được phân rã như Hình~\ref{fig:uc_hocbai}.

\begin{figure}[H]
\centering
\includegraphics[width=0.55\textwidth]{Hinhve/usecasephanra/usecasephanra_hocbaigiangtuongtac.png}
\caption{Biểu đồ use case phân rã Học bài giảng tương tác}
\label{fig:uc_hocbai}
\end{figure}

\emph{Xem video bài giảng} cho phép người học phát video; hệ thống tự động lưu lại tiến độ xem để người học có thể tiếp tục ở những lần sau. \emph{Trả lời hoạt động tương tác} là use case trung tâm: tại mỗi mốc thời gian được chèn hoạt động, người học phải trả lời câu hỏi (trắc nghiệm, điền khuyết, bài nghe, bài đọc hoặc bài viết tự luận) trước khi tiếp tục xem. \emph{Nộp bài và nhận chấm điểm tự động} chấm điểm câu trả lời và trả về phản hồi ngay lập tức.

\subsection{Biểu đồ use case phân rã Xem kết quả và phản hồi}
\label{subsection:uc_xemkq}
Use case \emph{Xem kết quả và phản hồi} của người học được phân rã như Hình~\ref{fig:uc_xemkq}.

\begin{figure}[H]
\centering
\includegraphics[width=0.55\textwidth]{Hinhve/usecasephanra/usecasephanra_xemketquavaphanhoi.png}
\caption{Biểu đồ use case phân rã Xem kết quả và phản hồi}
\label{fig:uc_xemkq}
\end{figure}

\emph{Xem điểm các bài đã làm} hiển thị điểm của từng bài giảng mà người học đã làm. \emph{Xem đáp án và nhận xét} cho phép người học xem lại từng câu trả lời cùng đáp án đúng và phản hồi của hệ thống. \emph{Theo dõi tiến độ khóa học} tổng hợp tiến độ và kết quả của người học trên toàn khóa học.

\subsection{Biểu đồ use case phân rã Quản lý khóa học, chương và bài giảng}
\label{subsection:uc_quanlykh}
Use case \emph{Quản lý khóa học, chương và bài giảng} của giảng viên được phân rã như Hình~\ref{fig:uc_quanlykh}.

\begin{figure}[H]
\centering
\includegraphics[width=0.55\textwidth]{Hinhve/usecasephanra/usecasephanra_quanlykhoahocchuongvabaigiang.png}
\caption{Biểu đồ use case phân rã Quản lý khóa học, chương và bài giảng}
\label{fig:uc_quanlykh}
\end{figure}

Ba use case con cho phép giảng viên tổ chức nội dung theo cấu trúc ba cấp. \emph{Quản lý khóa học} gồm tạo mới, cập nhật, đổi trạng thái và xóa khóa học. \emph{Quản lý chương} cho phép thêm, sửa, xóa và sắp xếp các chương trong khóa học. \emph{Quản lý bài giảng} cho phép thêm, sửa, xóa và sắp xếp các bài giảng trong từng chương; còn việc đưa video vào bài giảng do use case \emph{Tải lên video bài giảng} đảm nhận.

\subsection{Biểu đồ use case phân rã Phân tích video sinh nội dung tương tác}
\label{subsection:uc_phantich}
Use case \emph{Phân tích video sinh nội dung tương tác} của giảng viên được phân rã như Hình~\ref{fig:uc_phantich}.

\begin{figure}[H]
\centering
\includegraphics[width=0.55\textwidth]{Hinhve/usecasephanra/usecasephanra_phantichvideosinhnoidungtuongtac.png}
\caption{Biểu đồ use case phân rã Phân tích video sinh nội dung tương tác}
\label{fig:uc_phantich}
\end{figure}

\emph{Khởi tạo phân tích video} cho phép giảng viên bắt đầu quá trình xử lý tự động đối với video đã tải lên; khi đó hệ thống lần lượt trích xuất lời thoại bằng công nghệ nhận dạng tiếng nói, phân đoạn nội dung và dùng mô hình ngôn ngữ lớn để sinh các hoạt động tương tác bám sát từng đoạn. Vì quá trình này chạy nền và kéo dài, \emph{Theo dõi tiến độ phân tích} cho phép giảng viên xem trạng thái và tiến độ của tác vụ, còn \emph{Hủy tác vụ phân tích} cho phép dừng tác vụ khi cần thiết. Các bước trích xuất, phân đoạn và sinh nội dung do hệ thống thực hiện tự động sau khi được kích hoạt, được mô tả như một hoạt động trong quy trình nghiệp vụ ở Phần~\ref{subsection:quytrinh}.

\subsection{Biểu đồ use case phân rã Biên tập bản chép lời}
\label{subsection:uc_banchep}
Use case \emph{Biên tập bản chép lời} của giảng viên được phân rã như Hình~\ref{fig:uc_banchep}.

\begin{figure}[H]
\centering
\includegraphics[width=0.55\textwidth]{Hinhve/usecasephanra/usecasephanra_bientapbancheploi.png}
\caption{Biểu đồ use case phân rã Biên tập bản chép lời}
\label{fig:uc_banchep}
\end{figure}

\emph{Sửa nội dung lời thoại} cho phép giảng viên chỉnh sửa văn bản mà hệ thống trích xuất tự động nhằm khắc phục lỗi nhận dạng. \emph{Gộp hoặc tách phân đoạn} điều chỉnh cách bản chép lời được chia thành các đoạn nội dung. \emph{Điều chỉnh ranh giới phân đoạn} thay đổi mốc thời gian bắt đầu và kết thúc của mỗi đoạn, qua đó kiểm soát vị trí đặt hoạt động tương tác.

\subsection{Biểu đồ use case phân rã Biên tập hoạt động tương tác}
\label{subsection:uc_bientap}
Use case \emph{Biên tập hoạt động tương tác} của giảng viên được phân rã như Hình~\ref{fig:uc_bientap}.

\begin{figure}[H]
\centering
\includegraphics[width=0.5\textwidth]{Hinhve/usecasephanra/usecasephanra_bientaphoatdongtuongtac.png}
\caption{Biểu đồ use case phân rã Biên tập hoạt động tương tác}
\label{fig:uc_bientap}
\end{figure}

Sau khi hệ thống sinh nội dung tự động, giảng viên \emph{Xem trước hoạt động} để kiểm tra kết quả. Giảng viên có thể \emph{Thêm hoạt động thủ công} tại mốc thời gian mong muốn, \emph{Chỉnh sửa hoạt động} đã có hoặc \emph{Xóa hoạt động} không phù hợp. Khi chưa hài lòng, \emph{Sinh lại hoạt động bằng AI} tạo lại hoạt động cho một đoạn nội dung.

\subsection{Biểu đồ use case phân rã Quản lý thành viên và duyệt yêu cầu tham gia}
\label{subsection:uc_thanhvien}
Use case \emph{Quản lý thành viên và duyệt yêu cầu tham gia} của giảng viên được phân rã như Hình~\ref{fig:uc_thanhvien}.

\begin{figure}[H]
\centering
\includegraphics[width=0.55\textwidth]{Hinhve/usecasephanra/usecasephanra_quanlythanhvienvaduyetyeucauthamgia.png}
\caption{Biểu đồ use case phân rã Quản lý thành viên và duyệt yêu cầu tham gia}
\label{fig:uc_thanhvien}
\end{figure}

\emph{Xem danh sách thành viên} hiển thị những học viên hiện có trong khóa học. \emph{Thêm hoặc xóa thành viên} cho phép giảng viên trực tiếp thêm hoặc loại bỏ người học. \emph{Duyệt yêu cầu tham gia} cho phép giảng viên xem các yêu cầu đang chờ và chấp nhận hoặc từ chối từng yêu cầu.

\subsection{Biểu đồ use case phân rã Theo dõi kết quả học tập}
\label{subsection:uc_theodoi}
Use case \emph{Theo dõi kết quả học tập} của giảng viên được phân rã như Hình~\ref{fig:uc_theodoi}.

\begin{figure}[H]
\centering
\includegraphics[width=0.55\textwidth]{Hinhve/usecasephanra/usecasephanra_theodoiketquahoctap.png}
\caption{Biểu đồ use case phân rã Theo dõi kết quả học tập}
\label{fig:uc_theodoi}
\end{figure}

\emph{Xem thống kê kết quả khóa học} tổng hợp điểm và tỷ lệ hoàn thành của cả khóa học. \emph{Xem biểu đồ kết quả theo mốc thời gian} thể hiện điểm trung bình của người học tại từng đoạn của video. \emph{Phân tích kết quả từng câu hỏi} cho biết tỷ lệ trả lời đúng của mỗi hoạt động. \emph{Xem danh sách người học cần hỗ trợ} chỉ ra những người học có kết quả thấp để giảng viên kịp thời can thiệp.

\subsection{Biểu đồ use case phân rã Quản lý người dùng và phân quyền}
\label{subsection:uc_nguoidung}
Use case \emph{Quản lý người dùng và phân quyền} của quản trị viên được phân rã như Hình~\ref{fig:uc_nguoidung}.

\begin{figure}[H]
\centering
\includegraphics[width=0.55\textwidth]{Hinhve/usecasephanra/usecasephanra_quanlynguoidungvaphanquyen.png}
\caption{Biểu đồ use case phân rã Quản lý người dùng và phân quyền}
\label{fig:uc_nguoidung}
\end{figure}

\emph{Xem danh sách người dùng} liệt kê toàn bộ tài khoản trên hệ thống. \emph{Tạo tài khoản người dùng} cho phép quản trị viên thêm tài khoản mới kèm vai trò và trạng thái. \emph{Gán vai trò người dùng} thay đổi quyền của một tài khoản. \emph{Khóa hoặc mở khóa tài khoản} kiểm soát quyền truy cập của người dùng vào hệ thống.

\subsection{Biểu đồ use case phân rã Quản lý tổ chức và thành viên}
\label{subsection:uc_tochuc}
Use case \emph{Quản lý tổ chức và thành viên} của quản trị viên được phân rã như Hình~\ref{fig:uc_tochuc}.

\begin{figure}[H]
\centering
\includegraphics[width=0.55\textwidth]{Hinhve/usecasephanra/usecasephanra_quanlytochucvathanhvien.png}
\caption{Biểu đồ use case phân rã Quản lý tổ chức và thành viên}
\label{fig:uc_tochuc}
\end{figure}

\emph{Quản lý tổ chức} gồm tạo mới, cập nhật, đổi trạng thái và xóa tổ chức đào tạo. \emph{Xem danh sách thành viên tổ chức} hiển thị các tài khoản thuộc một tổ chức. \emph{Thêm hoặc xóa thành viên tổ chức} điều chỉnh thành phần của tổ chức. \emph{Gán vai trò thành viên} thiết lập vai trò (chủ sở hữu, quản trị, giảng viên, học viên) cho từng thành viên trong tổ chức.

\subsection{Biểu đồ use case phân rã Giám sát tác vụ phân tích}
\label{subsection:uc_giamsat}
Use case \emph{Giám sát tác vụ phân tích} của quản trị viên được phân rã như Hình~\ref{fig:uc_giamsat}.

\begin{figure}[H]
\centering
\includegraphics[width=0.55\textwidth]{Hinhve/usecasephanra/usecasephanra_giamsattacvuphantich.png}
\caption{Biểu đồ use case phân rã Giám sát tác vụ phân tích}
\label{fig:uc_giamsat}
\end{figure}

\emph{Xem danh sách tác vụ toàn hệ thống} cho phép quản trị viên theo dõi mọi tác vụ phân tích video đang chạy hoặc đã hoàn thành. \emph{Xem chi tiết và trạng thái tác vụ} hiển thị tiến độ, trạng thái và thông tin lỗi (nếu có) của một tác vụ. \emph{Hủy tác vụ} cho phép dừng một tác vụ đang chạy khi cần.

\subsection{Quy trình nghiệp vụ}
\label{subsection:quytrinh}
Quy trình nghiệp vụ quan trọng nhất của hệ thống là chuyển hóa một video bài giảng thụ động thành bài học có tương tác và tổ chức cho người học sử dụng. Đây không phải luồng sự kiện của một use case riêng lẻ mà là luồng hoạt động kết hợp nhiều use case của ba tác nhân Giảng viên, Hệ thống và Người học. Biểu đồ hoạt động mô tả quy trình này được thể hiện trong Hình~\ref{fig:quytrinh}.

\begin{figure}[H]
\centering
\includegraphics[width=\textwidth,height=0.92\textheight,keepaspectratio]{Hinhve/bieudohoatdong.png}
\caption{Biểu đồ hoạt động mô tả quy trình chuyển hóa video bài giảng thành bài học tương tác}
\label{fig:quytrinh}
\end{figure}

Quy trình bắt đầu khi giảng viên \emph{tạo bài giảng và tải video lên} hệ thống. Hệ thống tự động \emph{phân tích video và sinh các hoạt động tương tác}: trích xuất lời thoại, phân đoạn nội dung và dùng mô hình ngôn ngữ lớn để đề xuất hoạt động cho từng đoạn; nếu tác vụ phân tích thất bại, hệ thống báo lỗi để giảng viên chạy lại. Khi phân tích thành công, giảng viên \emph{xem trước và biên tập bản chép lời, phân đoạn và hoạt động}; nếu nội dung chưa đạt yêu cầu, giảng viên chỉnh sửa hoặc yêu cầu sinh lại rồi kiểm tra lại, ngược lại thì \emph{xuất bản khóa học}. Sau khi khóa học được xuất bản, người học \emph{gửi yêu cầu tham gia}; nếu giảng viên từ chối, hệ thống cập nhật trạng thái yêu cầu bị từ chối, còn nếu được duyệt thì người học \emph{học bài giảng tương tác} với mỗi câu trả lời được hệ thống \emph{chấm điểm và phản hồi tự động} ngay trong lúc xem. Cuối cùng, người học \emph{xem lại kết quả và phản hồi} của mình, còn giảng viên \emph{theo dõi kết quả học tập} để đánh giá mức độ tiếp thu và kịp thời hỗ trợ. Quy trình khép kín này cho thấy cách hệ thống giảm tải khâu soạn thảo thủ công cho giảng viên đồng thời mang lại trải nghiệm học tập chủ động, có phản hồi tức thời cho người học.
\section{Đặc tả chức năng}
\label{section:2.3}

\subsection{Đặc tả use case Tham gia khóa học}
\begin{table}[H]
\centering
\small
\renewcommand{\arraystretch}{1.3}
\begin{tabular}{|>{\bfseries\raggedright\arraybackslash}p{2.6cm}|p{11.2cm}|}
\hline
Tên use case & Tham gia khóa học \\ \hline
Tác nhân & Người học (tác nhân phụ: Giảng viên) \\ \hline
Mục đích & Cho phép người học gửi yêu cầu và được duyệt tham gia một khóa học để truy cập các bài giảng. \\ \hline
Tiền điều kiện & Người học đã đăng nhập và là thành viên của tổ chức chứa khóa học; khóa học đã được xuất bản và hiển thị trong danh sách. \\ \hline
Kịch bản chính & \begin{tabular}{@{}>{\raggedright\arraybackslash}p{0.9cm}>{\raggedright\arraybackslash}p{2.0cm}p{7.3cm}@{}}
\textbf{STT} & \textbf{Thực hiện bởi} & \textbf{Hành động} \\[2pt]
1 & Người học & Tìm kiếm và mở thông tin một khóa học đã xuất bản. \\[2pt]
2 & Người học & Gửi yêu cầu tham gia khóa học. \\[2pt]
3 & Hệ thống & Tạo yêu cầu ở trạng thái chờ duyệt và hiển thị cho giảng viên. \\[2pt]
4 & Giảng viên & Chấp nhận yêu cầu tham gia. \\[2pt]
5 & Hệ thống & Thêm người học làm thành viên khóa học và cập nhật trạng thái yêu cầu thành đã chấp nhận. \\[2pt]
\end{tabular} \\ \hline
Kịch bản khác (Luồng sự kiện thay thế) & \begin{tabular}{@{}>{\raggedright\arraybackslash}p{0.9cm}>{\raggedright\arraybackslash}p{2.0cm}p{7.3cm}@{}}
\textbf{STT} & \textbf{Thực hiện bởi} & \textbf{Hành động} \\[2pt]
4a & Giảng viên & Từ chối yêu cầu; hệ thống cập nhật trạng thái bị từ chối và không thêm người học vào khóa học. \\[2pt]
\end{tabular} \\ \hline
Hậu điều kiện & Người học trở thành thành viên khóa học và truy cập được các bài giảng trong đó. \\ \hline
\end{tabular}
\caption{Đặc tả use case Tham gia khóa học}
\label{table:uc_thamgia}
\end{table}

\subsection{Đặc tả use case Học bài giảng tương tác}
\begin{table}[H]
\centering
\small
\renewcommand{\arraystretch}{1.3}
\begin{tabular}{|>{\bfseries\raggedright\arraybackslash}p{2.6cm}|p{11.2cm}|}
\hline
Tên use case & Học bài giảng tương tác \\ \hline
Tác nhân & Người học \\ \hline
Mục đích & Người học vừa xem video bài giảng vừa trả lời các hoạt động tương tác xen kẽ và được chấm điểm, phản hồi tự động. \\ \hline
Tiền điều kiện & Người học là thành viên khóa học; bài giảng có video và hoạt động tương tác; chưa vượt quá số lần làm bài cho phép. \\ \hline
Kịch bản chính & \begin{tabular}{@{}>{\raggedright\arraybackslash}p{0.9cm}>{\raggedright\arraybackslash}p{2.0cm}p{7.3cm}@{}}
\textbf{STT} & \textbf{Thực hiện bởi} & \textbf{Hành động} \\[2pt]
1 & Người học & Mở bài giảng và bắt đầu xem video. \\[2pt]
2 & Hệ thống & Phát video; tại mỗi mốc thời gian có hoạt động thì tạm dừng và hiển thị hoạt động tương tác. \\[2pt]
3 & Người học & Trả lời hoạt động được hiển thị. \\[2pt]
4 & Hệ thống & Chấm điểm câu trả lời và hiển thị phản hồi tức thời. \\[2pt]
5 & Người học & Tiếp tục xem và trả lời cho đến hết, sau đó nộp bài. \\[2pt]
6 & Hệ thống & Lưu kết quả trả lời, điểm số và tiến độ xem; cập nhật tiến độ học tập của người học sau khi nộp bài. \\[2pt]
\end{tabular} \\ \hline
Kịch bản khác (Luồng sự kiện thay thế) & \begin{tabular}{@{}>{\raggedright\arraybackslash}p{0.9cm}>{\raggedright\arraybackslash}p{2.0cm}p{7.3cm}@{}}
\textbf{STT} & \textbf{Thực hiện bởi} & \textbf{Hành động} \\[2pt]
1a & Người học & Quay lại bài đang học dở; hệ thống khôi phục tiến độ xem đã lưu. \\[2pt]
4a & Hệ thống & Nếu chế độ phản hồi đặt là ẩn đáp án thì chỉ ghi nhận câu trả lời, chưa lộ đáp án đúng. \\[2pt]
\end{tabular} \\ \hline
Hậu điều kiện & Kết quả bài làm và tiến độ học của người học được lưu lại. \\ \hline
\end{tabular}
\caption{Đặc tả use case Học bài giảng tương tác}
\label{table:uc_hocbai}
\end{table}

\subsection{Đặc tả use case Phân tích video sinh nội dung tương tác}
\begin{table}[H]
\centering
\small
\renewcommand{\arraystretch}{1.3}
\begin{tabular}{|>{\bfseries\raggedright\arraybackslash}p{2.6cm}|p{11.2cm}|}
\hline
Tên use case & Phân tích video sinh nội dung tương tác \\ \hline
Tác nhân & Giảng viên \\ \hline
Mục đích & Tự động trích xuất nội dung video và sinh các hoạt động tương tác bằng trí tuệ nhân tạo. \\ \hline
Tiền điều kiện & Bài giảng đã có video được tải lên hệ thống. \\ \hline
Kịch bản chính & \begin{tabular}{@{}>{\raggedright\arraybackslash}p{0.9cm}>{\raggedright\arraybackslash}p{2.0cm}p{7.3cm}@{}}
\textbf{STT} & \textbf{Thực hiện bởi} & \textbf{Hành động} \\[2pt]
1 & Giảng viên & Khởi tạo tác vụ phân tích cho bài giảng. \\[2pt]
2 & Hệ thống & Tạo tác vụ, đưa vào hàng đợi và xử lý ở chế độ nền. \\[2pt]
3 & Hệ thống & Trích xuất lời thoại từ video bằng công nghệ nhận dạng tiếng nói. \\[2pt]
4 & Hệ thống & Phân đoạn nội dung và dùng mô hình ngôn ngữ lớn sinh hoạt động tương tác cho từng đoạn. \\[2pt]
5 & Hệ thống & Cập nhật trạng thái tác vụ thành hoàn thành và lưu các hoạt động đã sinh. \\[2pt]
6 & Giảng viên & Theo dõi tiến độ và xem kết quả khi tác vụ hoàn thành. \\[2pt]
\end{tabular} \\ \hline
Kịch bản khác (Luồng sự kiện thay thế) & \begin{tabular}{@{}>{\raggedright\arraybackslash}p{0.9cm}>{\raggedright\arraybackslash}p{2.0cm}p{7.3cm}@{}}
\textbf{STT} & \textbf{Thực hiện bởi} & \textbf{Hành động} \\[2pt]
4a & Hệ thống & Nếu xử lý gặp lỗi thì cập nhật trạng thái thất bại kèm thông tin lỗi; giảng viên có thể chạy lại tác vụ. \\[2pt]
*a & Giảng viên & Trong khi tác vụ đang chạy, hủy tác vụ; hệ thống dừng và cập nhật trạng thái đã hủy. \\[2pt]
\end{tabular} \\ \hline
Hậu điều kiện & Bài giảng có tập hoạt động tương tác do trí tuệ nhân tạo sinh, sẵn sàng để biên tập. \\ \hline
\end{tabular}
\caption{Đặc tả use case Phân tích video sinh nội dung tương tác}
\label{table:uc_phantich}
\end{table}

\subsection{Đặc tả use case Biên tập hoạt động tương tác}
\begin{table}[H]
\centering
\small
\renewcommand{\arraystretch}{1.3}
\begin{tabular}{|>{\bfseries\raggedright\arraybackslash}p{2.6cm}|p{11.2cm}|}
\hline
Tên use case & Biên tập hoạt động tương tác \\ \hline
Tác nhân & Giảng viên \\ \hline
Mục đích & Giảng viên kiểm duyệt và tinh chỉnh các hoạt động tương tác trước khi đưa tới người học. \\ \hline
Tiền điều kiện & Bài giảng đã có hoạt động tương tác (do AI sinh hoặc thêm thủ công). \\ \hline
Kịch bản chính & \begin{tabular}{@{}>{\raggedright\arraybackslash}p{0.9cm}>{\raggedright\arraybackslash}p{2.0cm}p{7.3cm}@{}}
\textbf{STT} & \textbf{Thực hiện bởi} & \textbf{Hành động} \\[2pt]
1 & Giảng viên & Mở danh sách hoạt động tương tác của bài giảng. \\[2pt]
2 & Hệ thống & Hiển thị danh sách các hoạt động theo thứ tự sắp đặt trong bài giảng. \\[2pt]
3 & Giảng viên & Chỉnh sửa nội dung một hoạt động. \\[2pt]
4 & Hệ thống & Lưu thay đổi. \\[2pt]
\end{tabular} \\ \hline
Kịch bản khác (Luồng sự kiện thay thế) & \begin{tabular}{@{}>{\raggedright\arraybackslash}p{0.9cm}>{\raggedright\arraybackslash}p{2.0cm}p{7.3cm}@{}}
\textbf{STT} & \textbf{Thực hiện bởi} & \textbf{Hành động} \\[2pt]
3a & Giảng viên & Thêm một hoạt động thủ công tại mốc thời gian mong muốn. \\[2pt]
3b & Giảng viên & Xóa một hoạt động không phù hợp. \\[2pt]
3c & Giảng viên & Yêu cầu sinh lại hoạt động bằng AI; hệ thống sinh lại và thay thế hoạt động đó. \\[2pt]
\end{tabular} \\ \hline
Hậu điều kiện & Tập hoạt động tương tác được cập nhật theo chỉnh sửa của giảng viên. \\ \hline
\end{tabular}
\caption{Đặc tả use case Biên tập hoạt động tương tác}
\label{table:uc_bientap}
\end{table}

\subsection{Đặc tả use case Theo dõi kết quả học tập}
\begin{table}[H]
\centering
\small
\renewcommand{\arraystretch}{1.3}
\begin{tabular}{|>{\bfseries\raggedright\arraybackslash}p{2.6cm}|p{11.2cm}|}
\hline
Tên use case & Theo dõi kết quả học tập \\ \hline
Tác nhân & Giảng viên \\ \hline
Mục đích & Giảng viên xem các thống kê về kết quả làm bài và mức độ tiếp thu của người học để đánh giá và hỗ trợ kịp thời. \\ \hline
Tiền điều kiện & Khóa học có người học đã thực hiện bài giảng. \\ \hline
Kịch bản chính & \begin{tabular}{@{}>{\raggedright\arraybackslash}p{0.9cm}>{\raggedright\arraybackslash}p{2.0cm}p{7.3cm}@{}}
\textbf{STT} & \textbf{Thực hiện bởi} & \textbf{Hành động} \\[2pt]
1 & Giảng viên & Mở trang thống kê kết quả của khóa học hoặc bài giảng. \\[2pt]
2 & Hệ thống & Hiển thị thống kê tổng hợp gồm điểm và tỷ lệ hoàn thành của người học. \\[2pt]
3 & Giảng viên & Xem biểu đồ kết quả làm bài theo từng mốc thời gian của video. \\[2pt]
4 & Hệ thống & Hiển thị tỷ lệ trả lời đúng theo từng câu hỏi và danh sách người học cần hỗ trợ. \\[2pt]
\end{tabular} \\ \hline
Kịch bản khác (Luồng sự kiện thay thế) & \begin{tabular}{@{}>{\raggedright\arraybackslash}p{0.9cm}>{\raggedright\arraybackslash}p{2.0cm}p{7.3cm}@{}}
\textbf{STT} & \textbf{Thực hiện bởi} & \textbf{Hành động} \\[2pt]
2a & Hệ thống & Nếu chưa có người học làm bài thì hiển thị thông báo chưa có dữ liệu. \\[2pt]
\end{tabular} \\ \hline
Hậu điều kiện & Giảng viên nắm được tình hình học tập của khóa học để điều chỉnh nội dung hoặc hỗ trợ người học. \\ \hline
\end{tabular}
\caption{Đặc tả use case Theo dõi kết quả học tập}
\label{table:uc_theodoi}
\end{table}


\section{Yêu cầu phi chức năng}
\label{section:2.4}
Bên cạnh các chức năng đã đặc tả, hệ thống cần đáp ứng một số yêu cầu phi chức năng để bảo đảm chất lượng vận hành và trải nghiệm sử dụng. Phần này trình bày các yêu cầu quan trọng nhất theo từng khía cạnh.

\subsection{Tính bảo mật - Security}
Hệ thống phục vụ nhiều nhóm người dùng với vai trò khác nhau gồm khách, người học, giảng viên và quản trị viên, đồng thời tổ chức dữ liệu theo từng tổ chức đào tạo và khóa học riêng biệt. Chính vì thế, hệ thống cần xác thực danh tính người dùng và phân quyền chặt chẽ để mỗi người chỉ truy cập được những chức năng và dữ liệu thuộc phạm vi vai trò của mình, không thể xem hay can thiệp vào dữ liệu của tổ chức hay khóa học khác. Ngoài ra, các yêu cầu truy cập chức năng và dữ liệu đều phải qua xác thực và kiểm tra quyền, mọi dữ liệu đầu vào phải được kiểm tra tính hợp lệ trước khi xử lý, còn mật khẩu người dùng phải được lưu dưới dạng băm thay vì lưu trực tiếp.

\subsection{Hiệu năng - Performance}
Phần lớn thao tác của người dùng như đăng nhập, mở khóa học, xem bài giảng hay nộp một hoạt động tương tác đều mang tính tức thời, nên hệ thống cần phản hồi các thao tác này trong khoảng dưới ba giây để không làm gián đoạn trải nghiệm học tập. Riêng việc phân tích video để sinh hoạt động tương tác là tác vụ nặng và kéo dài. Vì vậy, hệ thống xử lý tác vụ này ở chế độ nền thay vì bắt người dùng chờ đồng bộ. Sau khi gửi yêu cầu, giảng viên có thể tiếp tục thao tác khác và theo dõi tiến độ cho tới khi tác vụ hoàn tất. Cách tách biệt giữa các thao tác tương tác và những tác vụ phân tích chạy dài giúp giao diện luôn giữ được độ phản hồi tốt ngay cả khi hệ thống đang xử lý nhiều video cùng lúc.

\subsection{Khả năng mở rộng - Scalability}
Một tổ chức đào tạo có thể có nhiều khóa học với đông người học cùng truy cập, đặc biệt vào các đợt cao điểm như đầu kỳ học. Do đó, hệ thống cần phục vụ đồng thời người học của một vài lớp học, tương đương vài chục người dùng truy cập cùng lúc mà vẫn hoạt động ổn định, đồng thời xử lý song song nhiều tác vụ phân tích video thông qua một nhóm tiến trình xử lý có thể điều chỉnh số lượng theo tài nguyên máy chủ. Bên cạnh đó, khi kho video bài giảng ngày càng lớn và xuất hiện những video có thời lượng dài, hệ thống cần phân tích được các video kích thước lớn bằng cách chia nhỏ chúng thành từng đoạn để xử lý lần lượt thay vì nạp toàn bộ vào bộ nhớ.

\subsection{Độ tin cậy - Reliability}
Vì một tác vụ phân tích video có thể kéo dài và đi qua nhiều bước, hệ thống cần theo dõi được trạng thái của từng tác vụ và có khả năng phục hồi khi tác vụ bị gián đoạn do sự cố, thay vì phải thực hiện lại từ đầu. Đối với các lỗi tạm thời phát sinh từ những dịch vụ bên ngoài, hệ thống tự động thử lại trước khi báo lỗi; nếu tác vụ vẫn thất bại, hệ thống ghi nhận lỗi và cho phép giảng viên chạy lại. Ngoài ra, các thao tác ghi dữ liệu quan trọng như lưu kết quả sinh hoạt động hay lưu bài làm của người học phải bảo đảm tính toàn vẹn và không tạo ra dữ liệu trùng lặp ngay cả khi tác vụ được chạy lại nhiều lần.

\subsection{Tính dễ dùng - Usability}
Người dùng của hệ thống chủ yếu là giảng viên và người học, không yêu cầu kiến thức kỹ thuật chuyên sâu, nên giao diện cần trực quan và sử dụng hoàn toàn bằng tiếng Việt. Trong quá trình học, người học cần nhận được phản hồi tức thời sau mỗi hoạt động tương tác để biết kết quả và điều chỉnh cách học, còn tiến độ xem video được lưu lại tự động để có thể tiếp tục ở những lần truy cập sau. Đối với giảng viên, hệ thống cần hiển thị rõ tiến độ của các tác vụ phân tích đang chạy và cho phép xem trước và chỉnh sửa các hoạt động do trí tuệ nhân tạo sinh ra trước khi đưa tới người học.

\subsection{Tính dễ bảo trì - Maintainability}
Hệ thống gồm nhiều thành phần như giao diện web, dịch vụ xử lý phía máy chủ và các dịch vụ trí tuệ nhân tạo, nên cần được tổ chức theo kiến trúc tách lớp rõ ràng để thuận tiện cho việc phát triển và bảo trì. Giao diện trao đổi dữ liệu giữa máy chủ và giao diện người dùng nên được mô tả tập trung và sinh mã tự động cho cả hai phía từ một định nghĩa chung; tương tự, các truy vấn cơ sở dữ liệu cũng được sinh mã từ mô tả tập trung. Cách làm này giúp giảm sai sót và bảo đảm tính nhất quán mỗi khi hệ thống thay đổi. Nhờ đó, việc bổ sung chức năng mới hay điều chỉnh chức năng hiện có có thể thực hiện mà ít ảnh hưởng tới các phần còn lại của hệ thống.

Như vậy, Chương~\ref{chapter:Related_works} đã khảo sát hiện trạng, xác định các tác nhân và chức năng tổng quan, đặc tả chi tiết những chức năng quan trọng cùng các yêu cầu phi chức năng mà hệ thống cần đáp ứng. Đây là cơ sở để Chương~\ref{chapter:Methodology} trình bày nền tảng lý thuyết và lựa chọn công nghệ nhằm hiện thực hóa các yêu cầu đã đặt ra.


%%%%%%%%%%%%%%%%%%%%%%%%%%%%%%%%%%%

\end{document}
